% !TeX root = ../MQF_Manuale_Master.tex
%2multibyte Version: 5.50.0.2960 CodePage: 65001
% MQF_Manuale_Master_SW55_Memoir.tex
% Master del manuale MQF costruito sullo shell SW5.5 "Memoir Class for Configurable Typesetting".
% Versione modulare: i capitoli scritti sono inclusi con \input.%
% POSIZIONE DEI FILE
% Questo master deve stare nella cartella 01_Manuale.
% I capitoli devono stare nella sottocartella 01_Manuale/Capitoli.
% Non mettere il master dentro Capitoli.
% I capitoli futuri restano come traccia, ma non generano pagine nel documento.
% Evita pagine bianche dovute all'apertura dei capitoli su pagina dispari.
% ------------------------------------------------------------
% Nomi italiani
% ------------------------------------------------------------
% ------------------------------------------------------------
% Macro matematiche ufficiali del progetto MQF
% Coerenti con MQF_Notazione.tex
% ------------------------------------------------------------
% ------------------------------------------------------------
% Ambienti di tipo teorema.
% I nomi interni restano compatibili con lo shell SW.
% Le etichette stampate sono in italiano.
% ------------------------------------------------------------

%TCIDATA{OutputFilter=latex2.dll}
%TCIDATA{Version=5.50.0.2960}
%TCIDATA{LaTeXparent=1,1,MQF_Manuale_Master.tex}


% Capitoli/MQF_Cap_07_Python_Traiettorie_Pricing.tex

\chapter[Dalle traiettorie al prezzo]{Dalle traiettorie al prezzo: simulazione Monte Carlo in Python}

\label{cap:python-traiettorie-pricing}



I Capitoli 5 e 6 hanno costruito il quadro teorico necessario per
descrivere l'evoluzione dell'incertezza nel tempo. Il primo ha introdotto
processi stocastici, traiettorie, scenari e simulazione Monte Carlo; il
secondo ha sviluppato i principali modelli in tempo continuo, mostrando
come equazioni differenziali stocastiche, dipendenza tra fattori e
discretizzazione numerica consentano di rappresentare prezzi, tassi e
altre grandezze finanziarie dinamiche.

Questo capitolo compie il passaggio successivo. L'obiettivo non \`e
introdurre nuovi modelli stocastici, ma utilizzare quelli gi\`a
studiati per costruire una procedura completa di valutazione numerica.
Le traiettorie simulate cessano quindi di essere soltanto
rappresentazioni dell'evoluzione possibile di un fattore di rischio e
diventano l'input per determinare flussi di cassa, payoff e valori
attuali scenario per scenario.

Il problema applicativo assume cos\`i una struttura diversa da quella
dei capitoli teorici. Non basta specificare una dinamica continua e
discretizzarla correttamente. Occorre organizzare simultaneamente un
numero elevato di traiettorie, conservare la dipendenza tra i fattori,
ricostruire lungo ciascuno scenario le grandezze finanziarie rilevanti
e infine aggregare i valori ottenuti mediante una stima Monte Carlo.
La correttezza della valutazione dipende quindi sia dalla coerenza del
modello sia dalla qualit\`a della sua implementazione numerica.

Il filo conduttore del capitolo \`e costituito da un Caso Aula dedicato
alla valutazione di un'obbligazione indicizzata all'inflazione. Il valore
del titolo dipende congiuntamente dall'evoluzione futura dell'inflazione
e del tasso di interesse nominale. I due fattori sono modellizzati
mediante processi in tempo continuo e sono sottoposti a shock correlati.
Le loro traiettorie determinano, da un lato, l'indicizzazione delle
cedole e del rimborso e, dall'altro, i fattori di sconto utilizzati per
riportare i flussi al tempo iniziale.

Il caso consente di rendere operativi diversi elementi gi\`a introdotti
nel Capitolo 6: processi mean-reverting, dinamica CIR, shock gaussiani
correlati, discretizzazione su una griglia temporale e approssimazione di
integrali lungo la traiettoria. La novit\`a consiste nel collegare
questi elementi all'architettura finanziaria del contratto e nel
costruire, per ogni scenario simulato, un valore attuale complessivo.
La distribuzione dei valori scenario-specifici diventa quindi la base
per stimare il prezzo e valutarne la precisione numerica.

Particolare attenzione sar\`a dedicata ai controlli. Una simulazione
che produce risultati numerici senza errori di esecuzione non \`e
necessariamente corretta. Occorre verificare la struttura di
correlazione degli shock, la coerenza delle traiettorie con le
propriet\`a dei modelli utilizzati, la corretta costruzione dei flussi,
la decomposizione del valore attuale e la stabilit\`a della stima.
Saranno inoltre distinti due problemi differenti: l'incertezza dovuta
al numero finito di simulazioni e l'approssimazione introdotta dalla
discretizzazione temporale.

Il metodo di lavoro rimane quello introdotto nel Capitolo 4 per i
capitoli applicativi. La Scheda Caso definisce la specifica del problema,
il flusso logico-teorico precede la costruzione del notebook e le tappe
computazionali devono mantenere riconoscibili input, operazioni, output
e controlli. Tali principi non vengono quindi ripresentati in questo
capitolo: vengono utilizzati come struttura di lavoro gi\`a acquisita,
concentrando l'attenzione sugli aspetti specifici della simulazione di
processi in tempo continuo e del pricing per traiettorie.

Accanto al Caso Aula viene proposto un Caso TakeHome dedicato alla
valutazione Monte Carlo di una call asiatica su un indice azionario in
presenza di un tasso di interesse stocastico. Il contesto finanziario e
la struttura del payoff cambiano, ma rimane invariato il problema
fondamentale: trasformare fattori di rischio simulati in una quantit\`a
finanziaria dipendente dalla traiettoria e successivamente in una stima
di valore accompagnata da controlli numerici.

L'obiettivo complessivo del capitolo \`e quindi rendere operativo il
passaggio dalla formulazione di un modello stocastico alla valutazione
di un contratto finanziario. Python costituisce lo strumento con cui la
legge dei processi viene trasformata in un insieme di scenari
osservabili; il contratto trasforma tali scenari in flussi o payoff; la
simulazione Monte Carlo consente infine di aggregare i valori ottenuti
in una stima quantitativa verificabile.

\section{Dai processi stocastici alla valutazione per traiettorie}

Il problema affrontato in questo capitolo non richiede di costruire nuovamente
la teoria dei processi stocastici. Gli elementi necessari sono gi\`a stati
introdotti nei Capitoli 5 e 6 e devono ora essere utilizzati con una funzione
diversa: non soltanto descrivere l'evoluzione aleatoria di una grandezza, ma
costruire una valutazione finanziaria a partire dalle sue possibili
traiettorie.

La distinzione \`e importante. Nei capitoli precedenti l'attenzione era
rivolta principalmente alla specificazione del processo, alle sue propriet\`a
e alla relazione tra modello continuo e rappresentazione discreta. In una
procedura di pricing occorre compiere un ulteriore passaggio: associare a ogni
traiettoria simulata le grandezze economiche generate dal contratto e
trasformarle in un valore finanziario confrontabile tra scenari.

\subsection{Gli elementi gi\`a acquisiti}

Il Capitolo 5 ha introdotto la nozione di processo stocastico come famiglia di
variabili casuali indicizzate dal tempo e ha distinto il processo dalla
singola traiettoria osservata o simulata. Ha inoltre mostrato come la
simulazione Monte Carlo permetta di generare molte repliche indipendenti di
una stessa struttura probabilistica e di utilizzare il campione risultante
per stimare quantit\`a attese.

Indicando con
\[
\left\{X_{t_0}^{(s)},X_{t_1}^{(s)},\ldots,X_{t_N}^{(s)}\right\},
\qquad
s=1,\ldots,M,
\]
le traiettorie simulate di un processo sulla griglia
\[
0=t_0<t_1<\cdots<t_N=T,
\]
la simulazione produce quindi non un singolo esito, ma un insieme di
evoluzioni possibili dello stesso fattore di rischio.

Il Capitolo 6 ha fornito gli strumenti necessari per costruire tali
traiettorie quando il modello \`e formulato in tempo continuo. Una dinamica
del tipo
\[
dX_t
=
a(X_t,t)\,dt
+
b(X_t,t)\,dW_t
\]
pu\`o essere rappresentata numericamente su una griglia mediante uno schema
ricorsivo. Nel caso di Eulero--Maruyama,
\[
X_{t_{n+1}}
=
X_{t_n}
+
a(X_{t_n},t_n)\Delta t
+
b(X_{t_n},t_n)\sqrt{\Delta t}\,Z_{n+1},
\]
con
\[
Z_{n+1}\sim\mathcal{N}(0,1).
\]

Sono gi\`a stati inoltre introdotti i modelli che verranno utilizzati nel Caso
Aula, la costruzione di shock gaussiani correlati e la distinzione tra errore
di simulazione Monte Carlo ed errore dovuto alla discretizzazione temporale.
Questi elementi costituiscono quindi input teorici del capitolo corrente e
non devono essere nuovamente derivati.

\subsection{Dalla traiettoria simulata alla quantit\`a finanziaria}

Una traiettoria di un fattore di rischio non coincide ancora con un valore
finanziario. Per ottenere quest'ultimo occorre specificare come il contratto
utilizzi l'evoluzione simulata.

In alcuni problemi \`e sufficiente osservare il valore terminale
\[
X_T.
\]
In altri casi il risultato dipende invece da pi\`u date o dall'intera
traiettoria. Una media temporale, un indice cumulato, un fattore di sconto o
una sequenza di flussi possono dipendere da
\[
X_{t_0},X_{t_1},\ldots,X_{t_N}
\]
e non soltanto dal valore finale.

Questa distinzione modifica il problema computazionale. Per ogni scenario
\(s\), la simulazione deve conservare le informazioni necessarie per
costruire una quantit\`a del tipo
\[
G^{(s)}
=
G\left(
X_{t_0}^{(s)},
X_{t_1}^{(s)},
\ldots,
X_{t_N}^{(s)}
\right),
\]
dove \(G\) rappresenta la trasformazione finanziaria richiesta dal
contratto.

Nel pricing, tale trasformazione pu\`o comprendere pi\`u passaggi. Le
traiettorie dei fattori determinano i flussi o il payoff; i tassi simulati
determinano i fattori di sconto; dalla loro combinazione si ottiene quindi un
valore attuale specifico dello scenario,
\[
V^{(s)}.
\]

Solo dopo avere costruito
\[
V^{(1)},V^{(2)},\ldots,V^{(M)}
\]
ha senso applicare l'aggregazione Monte Carlo. La stima del valore assume
allora la forma
\[
\widehat V_M
=
\frac{1}{M}
\sum_{s=1}^{M}
V^{(s)}.
\]

Il punto essenziale \`e quindi che la simulazione non genera direttamente il
prezzo. Genera le traiettorie dei fattori di rischio; il contratto trasforma
tali traiettorie in valori scenario-specifici; la media Monte Carlo interviene
soltanto nell'ultima fase.

\subsection{La catena computazionale}

Il passaggio dai modelli studiati nei capitoli precedenti al pricing pu\`o
essere sintetizzato nella sequenza

\[
\begin{array}{c}
\text{modello stocastico}
\\[3pt]
\downarrow
\\[3pt]
\text{discretizzazione temporale}
\\[3pt]
\downarrow
\\[3pt]
\text{shock simulati}
\\[3pt]
\downarrow
\\[3pt]
\text{traiettorie dei fattori}
\\[3pt]
\downarrow
\\[3pt]
\text{flussi o payoff per scenario}
\\[3pt]
\downarrow
\\[3pt]
\text{valore attuale per scenario}
\\[3pt]
\downarrow
\\[3pt]
\text{stima Monte Carlo}.
\end{array}
\]

Questa catena costituisce la struttura computazionale specifica del capitolo.
Il metodo generale con cui il problema viene analizzato, scomposto e
controllato rimane invece quello introdotto nel Capitolo 4. Non vengono quindi
ripresentati la funzione della Scheda Caso, i regimi di utilizzo
dell'intelligenza artificiale o la struttura generale delle tappe del
notebook. Tali elementi vengono assunti come gi\`a acquisiti.

Ci\`o che deve essere chiarito ora \`e come realizzare concretamente la
catena precedente quando il numero delle traiettorie \`e elevato, i fattori di
rischio sono pi\`u di uno e le quantit\`a finanziarie dipendono dall'intero
percorso temporale. Il primo problema diventa quindi organizzare in Python una
simulazione nella quale scenari e date siano rappresentati in modo coerente e
possano essere utilizzati nelle successive operazioni di valutazione.